
\documentclass[11pt]{scrbook}
\usepackage[T1]{fontenc}

\usepackage[french]{babel}
% \usepackage{babelbib}

%General packages for basic document setup
\usepackage{amsmath,amsfonts, amsthm, graphicx}
\usepackage[hmargin=1.5cm,vmargin=3cm,]{geometry}

% Random text generator package
\usepackage{lipsum}

\usepackage{cite}
\usepackage{xurl}
% Hyperlink package
\usepackage{hyperref}
\hypersetup{
  colorlinks=true,    %active colors for links
  linkcolor=blue,     %internal link color
  urlcolor=blue,      %external link color
  citecolor=blue,
}
\usepackage{stmaryrd}

\usepackage{enumitem}
\setlist[enumerate,1]{label={\textup{(\roman*)}}}
\allowdisplaybreaks
% Verbatism package : used for code display 
\usepackage{fancyvrb}
\usepackage{minted} %Interpret code and display colors aswell.

% Header and footer package
\usepackage{fancyhdr, lastpage}
% Header and footer parameters
% Some basic elements are \thepage for the number of the current page,
% \pageref{LastPage} for the number of the last page (add * to remove 
% the hyperlink.
% \rightmark for the current section, \leftmark for the current chapter,
\pagestyle{fancy}


%Chapter-Style : Only use with documentclass that require chaters !
\usepackage[Glenn]{fncychap}
\ChNameVar{\bfseries\Large\sffamily}
\ChTitleVar{\bfseries\Large\rmfamily}

%Colors package
\usepackage{xcolor}


% package for better refs
\usepackage[nameinlink]{cleveref}
% Some parameters to recognize tcbtheorem boxes as theorems, propositions, ...
\crefformat{theoremenv}{#2Théorème~#1#3}
\Crefformat{theoremenv}{#2Théorème~#1#3}

\crefformat{propositionenv}{#2Proposition~#1#3}
\Crefformat{propositionenv}{#2Proposition~#1#3}

\crefformat{lemmaenv}{#2Lemme~#1#3}
\Crefformat{lemmaenv}{#2Lemme~#1#3}

\crefformat{corollaryenv}{#2Corollaire~#1#3}
\Crefformat{corollaryenv}{#2Corollaire~#1#3}

\crefformat{definitionenv}{#2Définition~#1#3}
\Crefformat{definitionenv}{#2Définition~#1#3}

\crefformat{exampleenv}{#2Exemple~#1#3}
\Crefformat{exampleenv}{#2Exemple~#1#3}

\crefformat{remarkenv}{#2Remarque~#1#3}
\Crefformat{remarkenv}{#2Remarque~#1#3}

\crefformat{applicationenv}{#2Application~#1#3}
\Crefformat{applicationenv}{#2Application~#1#3}
\newcommand{\creflastconjunction}{ et }

%Figure package
\usepackage{tikz}
\usepackage{tikz-cd}

%Boxes packages, for theorems, definitions, ...
%REQUIRES tikz in order to work !
\usepackage[most]{tcolorbox}
\tcbuselibrary{breakable}
% if only tcbset is used, it is not possible to make nested environments
% Hence the use of tcbsetforeverylayer
\tcbsetforeverylayer{enhanced jigsaw,breakable,fonttitle=\bfseries}

% ----------- New environments ----------

\usepackage{aliascnt}
% Theorems

\newtheoremstyle{theoremstyle}% name of the style to be used
  {\topsep}% measure of space to leave above the theorem. E.g.: 3pt
  {\topsep}% measure of space to leave below the theorem. E.g.: 3pt
  {\itshape}% name of font to use in the body of the theorem
  {0pt}% measure of space to indent
  {\bfseries}% name of head font
  {}% punctuation between head and body
  { }% space after theorem head; " " = normal interword space
  {
    \textcolor{black}
    {\thmname{#1} \thmnumber{#2}{\thmnote{#3}}}
  }

\theoremstyle{theoremstyle}
\newtheorem{theoremenv}{Théorème}[chapter]
\newenvironment{theorem}[1]{
  \begin{tcolorbox}[frame hidden, boxrule=0pt ,
    borderline west={2pt}{0pt}{black}, colback=white]{}
    \begin{theoremenv}{#1}
  }{
  \end{theoremenv}
\end{tcolorbox}
}

\newtheoremstyle{proofstyle}% name of the style to be used
  {\topsep}% measure of space to leave above the theorem. E.g.: 3pt
  {\topsep}% measure of space to leave below the theorem. E.g.: 3pt
  {}% name of font to use in the body of the theorem
  {0pt}% measure of space to indent
  {\bfseries}% name of head font
  {\textcolor{black}:}% punctuation between head and body
  { }% space after theorem head; " " = normal interword space
  {
    \textcolor{black}
    {\thmname{#1}}
  }

\theoremstyle{proofstyle}
\newtheorem{proofenv}{Preuve}
% simple proof environment
\renewenvironment{proof}{
  \vspace{-0.59cm}
  \begin{tcolorbox}[frame hidden, boxrule=0pt ,
    borderline west={2pt}{0pt}{black}, colback=white,
%     borderline north={2pt}{0pt}{red, opacity=0.05}
    ]{}
  \begin{proofenv}
  }{
  \end{proofenv}
  \end{tcolorbox}
}

\newenvironment{soloproof}{
  \begin{tcolorbox}[frame hidden, boxrule=0pt ,
    borderline west={2pt}{0pt}{black}, colback=white,
%     borderline north={2pt}{0pt}{red, opacity=0.05}
    ]{}
  \begin{proofenv}
  }{
  \end{proofenv}
  \end{tcolorbox}
}
% if needed to make indentations in a proof, it is better to then use the following
% environments.
\newenvironment{proofe}{
  \vspace{-0.59cm}
  \begin{tcolorbox}[frame hidden, boxrule=0pt,borderline west={2pt}{0pt}{black},
      , colback=white
    ]{}
}{
\end{tcolorbox}
}

% one left rule
\newenvironment{proofr}{
  \vspace{-0.59cm}
  \begin{tcolorbox}[frame hidden, boxrule=0pt,borderline west={2pt}{0pt}{black},
    borderline west={2pt}{5mm}{black}, colback=white, left=10mm
    ]{}
}{
\end{tcolorbox}
}
% two left rule
\newenvironment{proofrr}{
  \vspace{-0.59cm}
  \begin{tcolorbox}[frame hidden, boxrule=0pt,borderline west={2pt}{0pt}{black},
    borderline west={1pt}{10mm}{black!80!white}, borderline west={2pt}{5mm}{black}, colback=white, left=15mm
    ]{}
}{
\end{tcolorbox}
}

% be careful as nested tcolorbox cannot break, so use this with precation if it is in a nested env.
\newenvironment{leftrule}[1]{
  \begin{tcolorbox}[frame hidden, boxrule=0pt, borderline west={2pt}{0pt}{black}, 
    opacityback=0]{(#1)}
}
{
\end{tcolorbox}
}

%Propositions
\newaliascnt{propositionenv}{theoremenv}

\theoremstyle{theoremstyle}
\newtheorem{propositionenv}[propositionenv]{Proposition}
\newenvironment{proposition}[1]{
  \begin{tcolorbox}[frame hidden, boxrule=0pt ,
    borderline west={2pt}{0pt}{black}, colback=white]{}
    \begin{propositionenv}{#1}
  }{
  \end{propositionenv}
\end{tcolorbox}
}
\aliascntresetthe{propositionenv}

%Lemmas
\newaliascnt{lemmaenv}{theoremenv}

\theoremstyle{theoremstyle}
\newtheorem{lemmaenv}[lemmaenv]{Lemme}
\newenvironment{lemma}[1]{
  \begin{tcolorbox}[frame hidden, boxrule=0pt ,
    borderline west={2pt}{0pt}{black}, colback=white]{}
    \begin{lemmaenv}{#1}
  }{
  \end{lemmaenv}
\end{tcolorbox}
}
\aliascntresetthe{lemmaenv}

%Corollaries
\newaliascnt{corollaryenv}{theoremenv}

\theoremstyle{theoremstyle}
\newtheorem{corollaryenv}[corollaryenv]{Corollaire}
\newenvironment{corollary}[1]{
  \begin{tcolorbox}[frame hidden, boxrule=0pt ,
    borderline west={2pt}{0pt}{black}, colback=white]{}
    \begin{corollaryenv}{#1}
  }{
  \end{corollaryenv}
\end{tcolorbox}
}
\aliascntresetthe{corollaryenv}

%Definitions
\newaliascnt{definitionenv}{theoremenv}

\newtheoremstyle{definitionstyle}% name of the style to be used
  {\topsep}% measure of space to leave above the definition. E.g.: 3pt
  {\topsep}% measure of space to leave below the definition. E.g.: 3pt
  {}% name of font to use in the body of the definition
  {0pt}% measure of space to indent
  {\bfseries}% name of head font
  {}% punctuation between head and body
  { }% space after definition head; " " = normal interword space
  {
    \textcolor{black}
    {\thmname{#1} \thmnumber{#2}{\thmnote{#3}}}
  }
\theoremstyle{definitionstyle}

\newtheorem{definitionenv}[definitionenv]{Définition}
\newenvironment{definition}[1]{
  \begin{tcolorbox}[frame hidden, boxrule=0pt ,
    borderline west={2pt}{0pt}{black}, colback=white]{}
    \begin{definitionenv}{#1}
  }{
  \end{definitionenv}
\end{tcolorbox}
}
\aliascntresetthe{definitionenv}

%Examples
\newaliascnt{exampleenv}{theoremenv}

\theoremstyle{definitionstyle}
\newtheorem{exampleenv}[exampleenv]{Exemple}
\newenvironment{example}{
  \begin{tcolorbox}[frame hidden, boxrule=0pt ,
    borderline west={2pt}{0pt}{black}, colback=white]{}
    \begin{exampleenv}{}
  }{
  \end{exampleenv}
\end{tcolorbox}
}
\aliascntresetthe{exampleenv}

%Remarks
\newaliascnt{remarkenv}{theoremenv}

\theoremstyle{definitionstyle}
\newtheorem{remarkenv}[remarkenv]{Remarque}
\newenvironment{remark}{
  \begin{tcolorbox}[frame hidden, boxrule=0pt ,
    borderline west={2pt}{0pt}{black}, colback=white]{}
    \begin{remarkenv}{}
  }{
  \end{remarkenv}
\end{tcolorbox}
}
\aliascntresetthe{remarkenv}

\newaliascnt{applicationenv}{theoremenv}

\theoremstyle{definitionstyle}
\newtheorem{applicationenv}[applicationenv]{Application}
\newenvironment{application}[1]{
  \begin{tcolorbox}[frame hidden, boxrule=0pt ,
    borderline west={2pt}{0pt}{black}, colback=white]{}
    \begin{applicationenv}{#1}
  }{
  \end{applicationenv}
\end{tcolorbox}
}
\aliascntresetthe{applicationenv}

\tcbset{texercisestyle/.style={arc=0.5mm, colframe=blue!25!yellow!90!white,
colback=blue!25!yellow!5!white, coltitle=blue!25!yellow!40!black,
fonttitle=\small\sffamily\bfseries, fontupper=\small, fontlower=\small,
listing options={style=tcblatex,texcsstyle=*\color{red!40!black}},
}}

\newtcolorbox[auto counter,list inside=exam]{texercise}[2][]{%
texercisestyle,
listing file={solutions/texercise\thetcbcounter.tex},
label={exe:#2},
record={\string\processsol{solutions/texercise\thetcbcounter.tex}{#2}},
title={Exercice \thetcbcounter\hfill\mdseries Solution à la page \pageref{sol:#2}},
list text={Exercice avec solution à la page \pageref{sol:#2}},#1}



% --- Exercises and solutions environments ---- 
%Use this package for correct use of exercise and answer environment.
% \usepackage{exercise}
% 
% % Config of the exercise package
% \renewcounter{Exercise}[chapter] %counter for the exercises that resets on each chapter.
% \renewcounter{Answer}[chapter]
% 
% % Header for the exercises
% \def\ExerciseHeader{\noindent\ExerciseHeaderDifficulty\textcolor{brown}{\textbf{Exercice \ExerciseHeaderNB} : \ExerciseHeaderTitle} \\ }
% 
% % Header for the answers
% \def\AnswerHeader{\noindent\textcolor{brown}{\textbf{Solution de l'exercice \ExerciseHeaderNB}}\\ }
% 
% %Exercise environment
% \newenvironment{exo}[1]{
%   \stepcounter{Exercise}
%   \begin{tcolorbox}[frame hidden, boxrule=0pt ,
%     borderline west={2pt}{0pt}{brown}, colback=brown!10!white]{}
%   \begin{Exercise}[title={#1},label={\the\value{chapter}-\the\value{Exercise}}, number={\the\value{Exercise}},
%     counter={Exercise}]
%   }{
%   % Next line is to create a hyperref to the solution
%   % [Solution \refAnswer{\ExerciseLabel}]
%   \end{Exercise} 
% \end{tcolorbox}
% }
% %Solution environment
% \newenvironment{ans}{
%   \begin{tcolorbox}[frame hidden, boxrule=0pt ,
%     borderline west={2pt}{0pt}{brown!50!white}, colback=brown!5!white]{}
%   \begin{Answer}[ref=\ExerciseLabel]
% }{\end{Answer} 
% \end{tcolorbox}
% }

%%%%% USAGE ADVICE : ALWAYS USE SOLUTION ENVIRONMENT RIGHT AFTER AN EXERCISE ENVIRONMENT OTHERWISE,
%%%%% IT WILL BREAK THE NUMBERING SYSTEM.

% Line spacing package
\usepackage{setspace}
\onehalfspacing
\usepackage{flushend}
\raggedbottom
% package used for multicolumns 
\usepackage{multicol}

% package for arrows in multiple equations for explanations (and more...)
\usepackage{witharrows}

\usepackage[bottom]{footmisc}
\usepackage{footnote}

\sloppy
\usepackage[final, nopatch=footnote]{microtype} % the option is to avoid the 
% Warning : Unable to apply patch 'footnote' 


\usepackage{amssymb}
\usepackage{faktor}


\DeclareTOCStyleEntries[pagenumberwidth=40pt]{tocline}{part,chapter,section,subsection}
% ----------- New commands ------------ %

\newcommand*{\blank}{\hspace*{1pt}}
\newcommand*{\MnR}[0]{\mathscr{M}_{n}(\mathbb{R})}
\newcommand*{\GLnR}[0]{\mathbf{GL}_{n}(\mathbb{R})}
\newcommand*{\transpose}[1]{{}^t #1}
\newcommand*{\GL}[0]{\textup{GL}}
\newcommand{\Alt}{\mathfrak{A}}
\newcommand*{\Sym}[0]{\mathfrak{S}}
\newcommand{\Stab}[0]{\textup{Stab}}
\newcommand{\End}[0]{\textup{End}}
\newcommand{\Fix}[0]{\textup{Fix}}
\newcommand*{\rg}[0]{\textup{rg}}
\renewcommand*{\Im}[0]{\textup{Im}}
\newcommand*{\Ker}[0]{\textup{Ker}}
\newcommand*{\Coker}[0]{\textup{Coker}}
\newcommand*{\Tr}[0]{\textup{Tr}}
\newcommand*{\norm}[1]{\left\lVert #1 \right\rVert}
\newcommand*{\indicatrice}[0]{\mathds{1}}
\newcommand*{\Vect}[1]{\textup{Vect}\left\{#1\right\}}
\newcommand*{\id}{\textup{id}}
\newcommand*{\Tor}{\textup{Tor}}
\newcommand*{\ev}{\textup{ev}}
\newcommand{\curvearrow}[1]{\Arrow[tikz={font={\tiny\mdseries}, text width=2cm}]{#1}}
\newcommand*{\proofspace}{\begin{proofe} \end{proofe}}
